\chapter{Unit 1: Workshop Instruments, Components, Trainer Module \& Display Devices}

\section{Section A: Multiple Choice Practice Questions (60 Curated Problems)}

\noindent The following 60 multiple-choice questions reflect the official workshop evaluation bank, categorized by fundamental learning outcome and laboratory testing difficulty.\vspace{3mm}

\noindent\textbf{Q1. Which component is commonly used to limit the flow of electric current?} \hfill [Tag: Familiarization of electronic components | Easy]
\begin{enumerate}[label=(\alph*),topsep=2pt,itemsep=1pt,partopsep=0pt,parsep=0pt]
    \item {Resistor}
    \item {Transistor}
    \item {Diode}
    \item {Capacitor}
\end{enumerate}
\vspace{2mm}

\noindent\textbf{Q2. Which component stores electrical energy in an electric field?} \hfill [Tag: Familiarization of electronic components | Easy]
\begin{enumerate}[label=(\alph*),topsep=2pt,itemsep=1pt,partopsep=0pt,parsep=0pt]
    \item {Capacitor}
    \item {Diode}
    \item {Fuse}
    \item {Resistor}
\end{enumerate}
\vspace{2mm}

\noindent\textbf{Q3. What is the main function of a diode in a simple circuit?} \hfill [Tag: Familiarization of electronic components | Easy]
\begin{enumerate}[label=(\alph*),topsep=2pt,itemsep=1pt,partopsep=0pt,parsep=0pt]
    \item {To store current in two directions}
    \item {To produce sound from light}
    \item {To increase resistance in all directions}
    \item {To allow current in one direction}
\end{enumerate}
\vspace{2mm}

\noindent\textbf{Q4. What does LED stand for?} \hfill [Tag: Familiarization of electronic components | Easy]
\begin{enumerate}[label=(\alph*),topsep=2pt,itemsep=1pt,partopsep=0pt,parsep=0pt]
    \item {Light Energy Detector}
    \item {Low Energy Device}
    \item {Linear Electrical Display}
    \item {Light Emitting Diode}
\end{enumerate}
\vspace{2mm}

\noindent\textbf{Q5. Which component is commonly used as an electronic switch or amplifier?} \hfill [Tag: Familiarization of electronic components | Easy]
\begin{enumerate}[label=(\alph*),topsep=2pt,itemsep=1pt,partopsep=0pt,parsep=0pt]
    \item {Resistor}
    \item {Inductor}
    \item {Capacitor}
    \item {Transistor}
\end{enumerate}
\vspace{2mm}

\noindent\textbf{Q6. What is the purpose of a fuse in an electronic circuit?} \hfill [Tag: Familiarization of electronic components | Easy]
\begin{enumerate}[label=(\alph*),topsep=2pt,itemsep=1pt,partopsep=0pt,parsep=0pt]
    \item {To store electrical charge}
    \item {To display a binary number}
    \item {To increase the supply voltage}
    \item {To protect against excessive current}
\end{enumerate}
\vspace{2mm}

\noindent\textbf{Q7. Which component is designed to store energy in a magnetic field?} \hfill [Tag: Familiarization of electronic components | Easy]
\begin{enumerate}[label=(\alph*),topsep=2pt,itemsep=1pt,partopsep=0pt,parsep=0pt]
    \item {Inductor}
    \item {Diode}
    \item {LED}
    \item {Resistor}
\end{enumerate}
\vspace{2mm}

\noindent\textbf{Q8. What is an integrated circuit, commonly called an IC?} \hfill [Tag: Familiarization of electronic components | Easy]
\begin{enumerate}[label=(\alph*),topsep=2pt,itemsep=1pt,partopsep=0pt,parsep=0pt]
    \item {A group of circuits on one chip}
    \item {A single wire with insulation}
    \item {A device used only for measuring voltage}
    \item {A mechanical switch with many contacts}
\end{enumerate}
\vspace{2mm}

\noindent\textbf{Q9. What is the main purpose of a digital trainer module?} \hfill [Tag: Familiarization of the digital trainer module | Easy]
\begin{enumerate}[label=(\alph*),topsep=2pt,itemsep=1pt,partopsep=0pt,parsep=0pt]
    \item {To measure room temperature only}
    \item {To charge batteries automatically}
    \item {To convert mechanical motion into sound}
    \item {To test and build digital circuits}
\end{enumerate}
\vspace{2mm}

\noindent\textbf{Q10. Which source on a digital trainer is commonly used to provide logic 0 or logic 1 inputs?} \hfill [Tag: Familiarization of the digital trainer module | Easy]
\begin{enumerate}[label=(\alph*),topsep=2pt,itemsep=1pt,partopsep=0pt,parsep=0pt]
    \item {Transformer core}
    \item {Heat sink}
    \item {Speaker}
    \item {Logic switch}
\end{enumerate}
\vspace{2mm}

\noindent\textbf{Q11. Which indicator on a digital trainer commonly shows the state of a logic signal?} \hfill [Tag: Familiarization of the digital trainer module | Easy]
\begin{enumerate}[label=(\alph*),topsep=2pt,itemsep=1pt,partopsep=0pt,parsep=0pt]
    \item {LED indicator}
    \item {Fuse holder}
    \item {Power transformer}
    \item {Variable resistor}
\end{enumerate}
\vspace{2mm}

\noindent\textbf{Q12. What is the function of a regulated power supply in a digital trainer?} \hfill [Tag: Familiarization of the digital trainer module | Easy]
\begin{enumerate}[label=(\alph*),topsep=2pt,itemsep=1pt,partopsep=0pt,parsep=0pt]
    \item {To provide a stable DC voltage}
    \item {To produce mechanical rotation}
    \item {To store programs permanently}
    \item {To measure resistance directly}
\end{enumerate}
\vspace{2mm}

\noindent\textbf{Q13. Which device on a digital trainer can provide regularly repeating digital pulses?} \hfill [Tag: Familiarization of the digital trainer module | Easy]
\begin{enumerate}[label=(\alph*),topsep=2pt,itemsep=1pt,partopsep=0pt,parsep=0pt]
    \item {Fuse}
    \item {Potentiometer}
    \item {Clock generator}
    \item {Loudspeaker}
\end{enumerate}
\vspace{2mm}

\noindent\textbf{Q14. What does a common seven-segment display show?} \hfill [Tag: Familiarization of the digital trainer module | Easy]
\begin{enumerate}[label=(\alph*),topsep=2pt,itemsep=1pt,partopsep=0pt,parsep=0pt]
    \item {Only resistance values}
    \item {Only analog waveforms}
    \item {Magnetic field lines}
    \item {Decimal digits}
\end{enumerate}
\vspace{2mm}

\noindent\textbf{Q15. Which instrument can generate a sine wave for testing an electronic circuit?} \hfill [Tag: Familiarization of input sources and display devices related to electronic circuits | Easy]
\begin{enumerate}[label=(\alph*),topsep=2pt,itemsep=1pt,partopsep=0pt,parsep=0pt]
    \item {Function generator}
    \item {Digital multimeter}
    \item {Seven-segment display}
    \item {Loudspeaker}
\end{enumerate}
\vspace{2mm}

\noindent\textbf{Q16. Which source provides a constant electrical voltage for a circuit?} \hfill [Tag: Familiarization of input sources and display devices related to electronic circuits | Easy]
\begin{enumerate}[label=(\alph*),topsep=2pt,itemsep=1pt,partopsep=0pt,parsep=0pt]
    \item {Microphone}
    \item {LED display}
    \item {DC power supply}
    \item {Oscilloscope}
\end{enumerate}
\vspace{2mm}

\noindent\textbf{Q17. Which device converts sound waves into an electrical signal?} \hfill [Tag: Familiarization of input sources and display devices related to electronic circuits | Easy]
\begin{enumerate}[label=(\alph*),topsep=2pt,itemsep=1pt,partopsep=0pt,parsep=0pt]
    \item {Microphone}
    \item {Relay}
    \item {Speaker}
    \item {Lamp}
\end{enumerate}
\vspace{2mm}

\noindent\textbf{Q18. Which device converts an electrical signal into sound?} \hfill [Tag: Familiarization of input sources and display devices related to electronic circuits | Easy]
\begin{enumerate}[label=(\alph*),topsep=2pt,itemsep=1pt,partopsep=0pt,parsep=0pt]
    \item {Photodiode}
    \item {Microphone}
    \item {Speaker}
    \item {Resistor}
\end{enumerate}
\vspace{2mm}

\noindent\textbf{Q19. Which instrument is used to observe the waveform of an electrical signal?} \hfill [Tag: Familiarization of input sources and display devices related to electronic circuits | Easy]
\begin{enumerate}[label=(\alph*),topsep=2pt,itemsep=1pt,partopsep=0pt,parsep=0pt]
    \item {Function switch}
    \item {Soldering iron}
    \item {Oscilloscope}
    \item {Battery holder}
\end{enumerate}
\vspace{2mm}

\noindent\textbf{Q20. Which display device uses liquid crystals to show information?} \hfill [Tag: Familiarization of input sources and display devices related to electronic circuits | Easy]
\begin{enumerate}[label=(\alph*),topsep=2pt,itemsep=1pt,partopsep=0pt,parsep=0pt]
    \item {LCD}
    \item {Diode}
    \item {Transformer}
    \item {LED}
\end{enumerate}
\vspace{2mm}

\noindent\textbf{Q21. A resistor with color bands brown, black, and red has an approximate resistance of:} \hfill [Tag: Familiarization of electronic components | Medium]
\begin{enumerate}[label=(\alph*),topsep=2pt,itemsep=1pt,partopsep=0pt,parsep=0pt]
    \item {100 $\Omega$}
    \item {10 k $\Omega$}
    \item {100 k $\Omega$}
    \item {1 k $\Omega$}
\end{enumerate}
\vspace{2mm}

\noindent\textbf{Q22. A circuit requires a 5 V supply and a component must be limited to 10 mA . What resistor value is closest for a simple series connection when the component drop is negligible?} \hfill [Tag: Familiarization of electronic components | Medium]
\begin{enumerate}[label=(\alph*),topsep=2pt,itemsep=1pt,partopsep=0pt,parsep=0pt]
    \item {250 $\Omega$}
    \item {5 k $\Omega$}
    \item {50 $\Omega$}
    \item {500 $\Omega$}
\end{enumerate}
\vspace{2mm}

\noindent\textbf{Q23. Which component is most suitable for allowing current in one direction while blocking it in the opposite direction?} \hfill [Tag: Familiarization of electronic components | Medium]
\begin{enumerate}[label=(\alph*),topsep=2pt,itemsep=1pt,partopsep=0pt,parsep=0pt]
    \item {Resistor}
    \item {Transformer}
    \item {Diode}
    \item {Capacitor}
\end{enumerate}
\vspace{2mm}

\noindent\textbf{Q24. An electrolytic capacitor is connected directly across a DC supply. Which practice is essential?} \hfill [Tag: Familiarization of electronic components | Medium]
\begin{enumerate}[label=(\alph*),topsep=2pt,itemsep=1pt,partopsep=0pt,parsep=0pt]
    \item {Reverse its marked terminals}
    \item {Connect it without polarity}
    \item {Match its polarity correctly}
    \item {Short its terminals initially}
\end{enumerate}
\vspace{2mm}

\noindent\textbf{Q25. In a transistor used as an electronic switch, the transistor is intended to operate mainly in which two regions?} \hfill [Tag: Familiarization of electronic components | Medium]
\begin{enumerate}[label=(\alph*),topsep=2pt,itemsep=1pt,partopsep=0pt,parsep=0pt]
    \item {Leakage and thermal}
    \item {Active and breakdown}
    \item {Reverse and avalanche}
    \item {Cutoff and saturation}
\end{enumerate}
\vspace{2mm}

\noindent\textbf{Q26. A fuse is placed in series with the supply of an electronic circuit primarily to:} \hfill [Tag: Familiarization of electronic components | Medium]
\begin{enumerate}[label=(\alph*),topsep=2pt,itemsep=1pt,partopsep=0pt,parsep=0pt]
    \item {Increase the circuit voltage}
    \item {Protect against excessive current}
    \item {Limit normal signal frequency}
    \item {Improve the circuit power factor}
\end{enumerate}
\vspace{2mm}

\noindent\textbf{Q27. Which component is commonly used to vary resistance manually in order to adjust voltage or signal level?} \hfill [Tag: Familiarization of electronic components | Medium]
\begin{enumerate}[label=(\alph*),topsep=2pt,itemsep=1pt,partopsep=0pt,parsep=0pt]
    \item {Crystal oscillator}
    \item {Rectifier diode}
    \item {Potentiometer}
    \item {Photodiode}
\end{enumerate}
\vspace{2mm}

\noindent\textbf{Q28. Before connecting an IC on a digital trainer module, the safest first action is to:} \hfill [Tag: Familiarization of the digital trainer module | Medium]
\begin{enumerate}[label=(\alph*),topsep=2pt,itemsep=1pt,partopsep=0pt,parsep=0pt]
    \item {Switch off the module supply}
    \item {Apply the maximum supply voltage}
    \item {Connect every input to ground}
    \item {Short the output terminals}
\end{enumerate}
\vspace{2mm}

\noindent\textbf{Q29. A digital trainer's logic switch is connected to an IC input. What is its main purpose?} \hfill [Tag: Familiarization of the digital trainer module | Medium]
\begin{enumerate}[label=(\alph*),topsep=2pt,itemsep=1pt,partopsep=0pt,parsep=0pt]
    \item {To measure analog resistance}
    \item {To provide selectable logic levels}
    \item {To amplify output current continuously}
    \item {To display the circuit waveform}
\end{enumerate}
\vspace{2mm}

\noindent\textbf{Q30. If a trainer module's logic probe indicates a high state at an output, the output voltage is generally:} \hfill [Tag: Familiarization of the digital trainer module | Medium]
\begin{enumerate}[label=(\alph*),topsep=2pt,itemsep=1pt,partopsep=0pt,parsep=0pt]
    \item {Near the circuit ground level}
    \item {Near the designated logic supply level}
    \item {Always equal to the input voltage}
    \item {Always equal to zero volts}
\end{enumerate}
\vspace{2mm}

\noindent\textbf{Q31. A seven-segment display on a trainer module shows an incorrect digit even though the logic inputs appear correct. Which check is most relevant first?} \hfill [Tag: Familiarization of the digital trainer module | Medium]
\begin{enumerate}[label=(\alph*),topsep=2pt,itemsep=1pt,partopsep=0pt,parsep=0pt]
    \item {Replace the display with a resistor}
    \item {Increase the supply without measuring}
    \item {Check the segment wiring and common terminal}
    \item {Disconnect the ground connection}
\end{enumerate}
\vspace{2mm}

\noindent\textbf{Q32. What is the main function of a clock generator on a digital trainer module?} \hfill [Tag: Familiarization of the digital trainer module | Medium]
\begin{enumerate}[label=(\alph*),topsep=2pt,itemsep=1pt,partopsep=0pt,parsep=0pt]
    \item {Convert DC directly into sound}
    \item {Measure the resistance of an IC}
    \item {Provide periodic digital pulses}
    \item {Regulate the brightness of all LEDs}
\end{enumerate}
\vspace{2mm}

\noindent\textbf{Q33. A digital IC input is left unconnected on a trainer module. Why can this cause unreliable operation?} \hfill [Tag: Familiarization of the digital trainer module | Medium]
\begin{enumerate}[label=(\alph*),topsep=2pt,itemsep=1pt,partopsep=0pt,parsep=0pt]
    \item {The input automatically becomes a regulated output}
    \item {The input permanently disables the power supply}
    \item {The input may assume an undefined logic state}
    \item {The input always produces a precise logic zero}
\end{enumerate}
\vspace{2mm}

\noindent\textbf{Q34. When testing a logic circuit on a trainer, why should the circuit ground and instrument ground be connected properly?} \hfill [Tag: Familiarization of the digital trainer module | Medium]
\begin{enumerate}[label=(\alph*),topsep=2pt,itemsep=1pt,partopsep=0pt,parsep=0pt]
    \item {To increase the logic supply frequency}
    \item {To convert binary signals into analog signals}
    \item {To prevent all output transitions}
    \item {To establish a common voltage reference}
\end{enumerate}
\vspace{2mm}

\noindent\textbf{Q35. A function generator is connected to an amplifier input. Which setting determines how large the applied signal is?} \hfill [Tag: Familiarization of input sources and display devices related to electronic circuits | Medium]
\begin{enumerate}[label=(\alph*),topsep=2pt,itemsep=1pt,partopsep=0pt,parsep=0pt]
    \item {Time-base scale}
    \item {Duty-cycle label}
    \item {Frequency}
    \item {Amplitude}
\end{enumerate}
\vspace{2mm}

\noindent\textbf{Q36. A sine wave displayed on an oscilloscope occupies four vertical divisions. If the vertical scale is 0.5 V/div , what is the peak-to-peak voltage?} \hfill [Tag: Familiarization of input sources and display devices related to electronic circuits | Medium]
\begin{enumerate}[label=(\alph*),topsep=2pt,itemsep=1pt,partopsep=0pt,parsep=0pt]
    \item {0.5 V}
    \item {4.0 V}
    \item {2.0 V}
    \item {1.0 V}
\end{enumerate}
\vspace{2mm}

\noindent\textbf{Q37. An oscilloscope shows five complete cycles across ten horizontal divisions. If the time scale is 0.2 ms/div , what is the approximate frequency?} \hfill [Tag: Familiarization of input sources and display devices related to electronic circuits | Medium]
\begin{enumerate}[label=(\alph*),topsep=2pt,itemsep=1pt,partopsep=0pt,parsep=0pt]
    \item {100 Hz}
    \item {25 kHz}
    \item {2.5 kHz}
    \item {250 Hz}
\end{enumerate}
\vspace{2mm}

\noindent\textbf{Q38. Which input device changes its resistance significantly in response to incident light and can be used in a light-sensing circuit?} \hfill [Tag: Familiarization of input sources and display devices related to electronic circuits | Medium]
\begin{enumerate}[label=(\alph*),topsep=2pt,itemsep=1pt,partopsep=0pt,parsep=0pt]
    \item {Piezoelectric buzzer}
    \item {LDR}
    \item {Thermistor}
    \item {Relay coil}
\end{enumerate}
\vspace{2mm}

\noindent\textbf{Q39. A microphone is connected to an electronic circuit as an input source. It primarily converts:} \hfill [Tag: Familiarization of input sources and display devices related to electronic circuits | Medium]
\begin{enumerate}[label=(\alph*),topsep=2pt,itemsep=1pt,partopsep=0pt,parsep=0pt]
    \item {Sound into an electrical signal}
    \item {Light into resistance}
    \item {DC voltage into magnetic flux}
    \item {Digital data into heat}
\end{enumerate}
\vspace{2mm}

\noindent\textbf{Q40. Which display device is most appropriate when a circuit must indicate the ON or OFF state of a digital output immediately?} \hfill [Tag: Familiarization of input sources and display devices related to electronic circuits | Medium]
\begin{enumerate}[label=(\alph*),topsep=2pt,itemsep=1pt,partopsep=0pt,parsep=0pt]
    \item {Thermistor}
    \item {Transformer}
    \item {LED}
    \item {Variable capacitor}
\end{enumerate}
\vspace{2mm}

\noindent\textbf{Q41. A resistor has the color bands brown, black, orange, and gold. During measurement with a multimeter, the value appears close to 9.8 k $\Omega$ . Which conclusion is most appropriate?} \hfill [Tag: Familiarization of electronic components | Hard]
\begin{enumerate}[label=(\alph*),topsep=2pt,itemsep=1pt,partopsep=0pt,parsep=0pt]
    \item {The gold band represents a temperature coefficient}
    \item {The resistor should measure near 100 $\Omega$}
    \item {The resistor is within its stated tolerance}
    \item {The resistor must be connected with reverse polarity}
\end{enumerate}
\vspace{2mm}

\noindent\textbf{Q42. A capacitor marked 104 J is connected to a circuit. Which interpretation is correct?} \hfill [Tag: Familiarization of electronic components | Hard]
\begin{enumerate}[label=(\alph*),topsep=2pt,itemsep=1pt,partopsep=0pt,parsep=0pt]
    \item {10.4 $\mu$ F with $\pm$ 5\% tolerance}
    \item {104 pF with $\pm$ 10\% tolerance}
    \item {100 nF with $\pm$ 5\% tolerance}
    \item {1 nF with $\pm$ 20\% tolerance}
\end{enumerate}
\vspace{2mm}

\noindent\textbf{Q43. A polarized electrolytic capacitor is installed across a DC supply with its negative terminal connected to the positive rail. What is the most likely consequence?} \hfill [Tag: Familiarization of electronic components | Hard]
\begin{enumerate}[label=(\alph*),topsep=2pt,itemsep=1pt,partopsep=0pt,parsep=0pt]
    \item {A permanent increase in dielectric strength}
    \item {Increased leakage and possible capacitor failure}
    \item {Reduced capacitance with no other effect}
    \item {Improved ripple filtering at higher frequency}
\end{enumerate}
\vspace{2mm}

\noindent\textbf{Q44. An LED and a 330 $\Omega$ resistor are connected in series to a 5 V supply. Assuming an LED forward drop of 2 V , which current is closest to the expected value?} \hfill [Tag: Familiarization of electronic components | Hard]
\begin{enumerate}[label=(\alph*),topsep=2pt,itemsep=1pt,partopsep=0pt,parsep=0pt]
    \item {9.1 mA}
    \item {66.0 mA}
    \item {6.1 mA}
    \item {15.2 mA}
\end{enumerate}
\vspace{2mm}

\noindent\textbf{Q45. A diode is tested using a digital multimeter in diode mode. It reads approximately 0.68 V in one direction and displays an open-circuit indication in the other. What does this most strongly indicate?} \hfill [Tag: Familiarization of electronic components | Hard]
\begin{enumerate}[label=(\alph*),topsep=2pt,itemsep=1pt,partopsep=0pt,parsep=0pt]
    \item {A reverse-biased germanium diode}
    \item {A normally functioning silicon diode}
    \item {An open diode junction in both directions}
    \item {A shorted silicon diode junction}
\end{enumerate}
\vspace{2mm}

\noindent\textbf{Q46. A transistor's data sheet identifies its terminals as E , B , and C , but the physical package orientation is unknown. Which method is safest for identifying the terminals before circuit connection?} \hfill [Tag: Familiarization of electronic components | Hard]
\begin{enumerate}[label=(\alph*),topsep=2pt,itemsep=1pt,partopsep=0pt,parsep=0pt]
    \item {Use resistance mode between every pair only}
    \item {Connect the supply and identify terminals by heating}
    \item {Assume the center lead is always the base}
    \item {Use the data sheet and verify junctions with diode tests}
\end{enumerate}
\vspace{2mm}

\noindent\textbf{Q47. A 5.1 V Zener diode is connected in reverse bias with a suitable series resistor. The supply is increased from 6 V to 12 V . Which behavior is expected if the current remains within its rated limits?} \hfill [Tag: Familiarization of electronic components | Hard]
\begin{enumerate}[label=(\alph*),topsep=2pt,itemsep=1pt,partopsep=0pt,parsep=0pt]
    \item {The Zener voltage becomes approximately twice as large}
    \item {The Zener current remains exactly unchanged}
    \item {The Zener voltage remains approximately constant}
    \item {The Zener behaves as a forward-biased LED}
\end{enumerate}
\vspace{2mm}

\noindent\textbf{Q48. A pushbutton input to a digital circuit changes unpredictably between logic states when pressed. Which component arrangement most directly addresses the problem?} \hfill [Tag: Familiarization of electronic components | Hard]
\begin{enumerate}[label=(\alph*),topsep=2pt,itemsep=1pt,partopsep=0pt,parsep=0pt]
    \item {A fuse inserted in series with the input}
    \item {A larger LED connected across the input}
    \item {A reverse-biased diode placed across the supply}
    \item {A pull-up or pull-down resistor with switch debouncing}
\end{enumerate}
\vspace{2mm}

\noindent\textbf{Q49. Two resistors are marked 1 k $\Omega$ and 2 k $\Omega$ . Which connection produces an equivalent resistance of approximately 667 $\Omega$ ?} \hfill [Tag: Familiarization of electronic components | Hard]
\begin{enumerate}[label=(\alph*),topsep=2pt,itemsep=1pt,partopsep=0pt,parsep=0pt]
    \item {Connection with one resistor open}
    \item {Series connection}
    \item {Connection across the same wire}
    \item {Parallel connection}
\end{enumerate}
\vspace{2mm}

\noindent\textbf{Q50. Before wiring a digital trainer circuit, the supply is set to the required voltage. Which additional practice best reduces the risk of damaging an integrated circuit?} \hfill [Tag: Familiarization of the digital trainer module | Hard]
\begin{enumerate}[label=(\alph*),topsep=2pt,itemsep=1pt,partopsep=0pt,parsep=0pt]
    \item {Increase the supply voltage to improve noise immunity}
    \item {Verify polarity and current-limit settings before insertion}
    \item {Short unused IC inputs to the positive rail}
    \item {Insert the IC while the trainer output is enabled}
\end{enumerate}
\vspace{2mm}

\noindent\textbf{Q51. A logic trainer's LED output indicator remains off when the measured output voltage is 3.8 V . Which diagnosis is most reasonable?} \hfill [Tag: Familiarization of the digital trainer module | Hard]
\begin{enumerate}[label=(\alph*),topsep=2pt,itemsep=1pt,partopsep=0pt,parsep=0pt]
    \item {The output voltage proves the LED is open-circuit}
    \item {The IC must be operating without any supply voltage}
    \item {The logic output is definitely at binary zero}
    \item {The indicator may be active-low or incorrectly referenced}
\end{enumerate}
\vspace{2mm}

\noindent\textbf{Q52. A digital trainer provides a clock generator and a manual pulse switch. A counter advances several times for one manual press. What is the most likely cause?} \hfill [Tag: Familiarization of the digital trainer module | Hard]
\begin{enumerate}[label=(\alph*),topsep=2pt,itemsep=1pt,partopsep=0pt,parsep=0pt]
    \item {The display decoder has changed the counter state}
    \item {Mechanical switch bounce creates multiple clock edges}
    \item {The counter has converted binary data into analog data}
    \item {The supply voltage is necessarily below the logic threshold}
\end{enumerate}
\vspace{2mm}

\noindent\textbf{Q53. A trainer module contains several unused inputs on a logic gate IC. Why should these inputs generally be assigned a defined logic level?} \hfill [Tag: Familiarization of the digital trainer module | Hard]
\begin{enumerate}[label=(\alph*),topsep=2pt,itemsep=1pt,partopsep=0pt,parsep=0pt]
    \item {Unused inputs automatically increase the output drive current}
    \item {Defined inputs convert the gate into an analog amplifier}
    \item {Floating inputs can respond to noise and cause unstable operation}
    \item {Floating inputs always force the output permanently low}
\end{enumerate}
\vspace{2mm}

\noindent\textbf{Q54. A trainer uses a common-ground digital supply. An oscilloscope probe is connected with its ground clip to a signal point rather than circuit ground. What is the primary risk?} \hfill [Tag: Familiarization of the digital trainer module | Hard]
\begin{enumerate}[label=(\alph*),topsep=2pt,itemsep=1pt,partopsep=0pt,parsep=0pt]
    \item {The probe will convert the signal into a differential measurement}
    \item {The circuit will become immune to ground noise}
    \item {The probe ground may short the signal to earth or circuit ground}
    \item {The displayed waveform will automatically double in frequency}
\end{enumerate}
\vspace{2mm}

\noindent\textbf{Q55. A logic trainer's regulated output is correct with no load but drops significantly when an IC circuit is connected. Which check should be performed first?} \hfill [Tag: Familiarization of the digital trainer module | Hard]
\begin{enumerate}[label=(\alph*),topsep=2pt,itemsep=1pt,partopsep=0pt,parsep=0pt]
    \item {Disconnect the ground to remove loading}
    \item {Replace every resistor with a lower resistance}
    \item {Increase the regulated voltage above the IC rating}
    \item {Inspect for wiring errors or an unintended short circuit}
\end{enumerate}
\vspace{2mm}

\noindent\textbf{Q56. A function generator is set to produce a 4 V pp sine wave with a 1 V DC offset. What are the maximum and minimum voltages relative to ground?} \hfill [Tag: Familiarization of input sources and display devices related to electronic circuits | Hard]
\begin{enumerate}[label=(\alph*),topsep=2pt,itemsep=1pt,partopsep=0pt,parsep=0pt]
    \item {+ 2 V and - 2 V}
    \item {+ 5 V and - 3 V}
    \item {+ 3 V and - 1 V}
    \item {+ 4 V and 0 V}
\end{enumerate}
\vspace{2mm}

\noindent\textbf{Q57. A function generator has a 50 $\Omega$ output and is connected to a high-impedance oscilloscope input. The generator display is calibrated for a 50 $\Omega$ load. What reading will commonly appear on the oscilloscope?} \hfill [Tag: Familiarization of input sources and display devices related to electronic circuits | Hard]
\begin{enumerate}[label=(\alph*),topsep=2pt,itemsep=1pt,partopsep=0pt,parsep=0pt]
    \item {Approximately twice the displayed amplitude}
    \item {Approximately the same amplitude with reversed polarity}
    \item {Approximately half the displayed amplitude}
    \item {Approximately one-quarter the displayed amplitude}
\end{enumerate}
\vspace{2mm}

\noindent\textbf{Q58. An oscilloscope displays a stable waveform only when its trigger level is set near the waveform's midpoint. What does this setting primarily control?} \hfill [Tag: Familiarization of input sources and display devices related to electronic circuits | Hard]
\begin{enumerate}[label=(\alph*),topsep=2pt,itemsep=1pt,partopsep=0pt,parsep=0pt]
    \item {The input impedance selected by the probe}
    \item {The vertical scale used to measure waveform amplitude}
    \item {The frequency generated by the signal source}
    \item {The voltage threshold at which sweep acquisition begins}
\end{enumerate}
\vspace{2mm}

\noindent\textbf{Q59. A seven-segment display is connected to a decoder whose outputs are active-low. The display segments illuminate only when their cathodes are driven low. Which display type is compatible?} \hfill [Tag: Familiarization of input sources and display devices related to electronic circuits | Hard]
\begin{enumerate}[label=(\alph*),topsep=2pt,itemsep=1pt,partopsep=0pt,parsep=0pt]
    \item {Differential-output display}
    \item {Common-cathode display}
    \item {Common-anode display}
    \item {Open-collector sensor display}
\end{enumerate}
\vspace{2mm}

\noindent\textbf{Q60. A digital circuit drives an LED directly from a logic output. The LED is brighter than expected and the IC becomes warm. Which change is most appropriate?} \hfill [Tag: Familiarization of input sources and display devices related to electronic circuits | Hard]
\begin{enumerate}[label=(\alph*),topsep=2pt,itemsep=1pt,partopsep=0pt,parsep=0pt]
    \item {Remove the LED's polarity check from the circuit}
    \item {Connect the LED directly across the supply rails}
    \item {Add a correctly sized series current-limiting resistor}
    \item {Increase the logic output voltage above its rating}
\end{enumerate}
\vspace{2mm}


\newpage
\section{Section B: Comprehensive Subjective Problem with Detailed Analytical Rubric}

\begin{subjectiveq}{1}{1}
A carbon-film resistor carries four color bands in order: Yellow, Violet, Orange, and Gold. Determine its nominal resistance, upper and lower tolerance limits, and explain why measuring this resistor while still connected inside a powered-on circuit will lead to erroneous DMM readings.
\end{subjectiveq}

\begin{solbox}
\textbf{1. Color Code Decoding:}\\
Yellow corresponds to digit 4, Violet to digit 7, and Orange corresponds to multiplier $10^3 = 1000\,\Omega$. The fourth band (Gold) represents a tolerance of $\pm 5\%$.\\
Nominal Resistance: $R = 47 \times 10^3\,\Omega = 47\,\mathrm{k}\Omega$.\\
Tolerance range: $\pm 5\% \times 47\,\mathrm{k}\Omega = \pm 2.35\,\mathrm{k}\Omega$.\\
Upper Limit: $R_{\max} = 47 + 2.35 = 49.35\,\mathrm{k}\Omega$.\\
Lower Limit: $R_{\min} = 47 - 2.35 = 44.65\,\mathrm{k}\Omega$.\\\\
\textbf{2. In-Circuit Measurement Error Mechanism:}\\
When a DMM is set to resistance mode, it injects an internal calibrated test current $I_{\mathrm{test}}$ and measures the resultant terminal voltage $V = I_{\mathrm{test}} R$. If the circuit is powered on or if parallel circuit paths exist:\\
(a) External circuit potentials corrupt the DMM's high-impedance voltage measurement ADC, potentially destroying the meter input fuse.\\
(b) Any parallel branches shunt a fraction of the test current away from the target resistor, yielding an equivalent parallel resistance $R_{\mathrm{meas}} = R \parallel R_{\mathrm{parallel}} < R$. Hence, resistors must always be tested de-energized and with at least one leg lifted from the breadboard.
\end{solbox}

\begin{examinertip}
Examiner Rubric (10 Marks):\\
- Correct decoding of color bands to digits \& multiplier: 3 Marks\\
- Correct nominal value and $\pm 5\%$ tolerance band calculation: 2 Marks\\
- Detailed physical explanation of DMM constant-current ohmmeter circuit: 3 Marks\\
- Discussion of parallel loading paths and breadboard isolation rule: 2 Marks
\end{examinertip}


\section{Section C: Laboratory Viva-Voce Questions with Authoritative Solutions}

\noindent\textbf{Viva Question:} What is the primary difference between passive and active electronic components?\\
\textbf{Technical Answer:} Passive components (resistors, capacitors, inductors) cannot amplify electrical signal power and do not require external power supplies to function. Active components (diodes, bipolar transistors, operational amplifiers) can control electron flow and amplify power, requiring an external DC operating bias.
\vspace{3mm}

\noindent\textbf{Viva Question:} Why does a Digital Storage Oscilloscope (DSO) require high input impedance (typically 1 M$\Omega$)?\\
\textbf{Technical Answer:} A 1 M$\Omega$ input impedance minimizes the loading effect on the circuit under test, preventing the oscilloscope probe from drawing significant signal current and perturbing the observed waveform amplitude or timing.
\vspace{3mm}

\noindent\textbf{Viva Question:} What is the purpose of the current-limiting resistor connected in series with a Light Emitting Diode (LED)?\\
\textbf{Technical Answer:} An LED exhibits an exponential forward V-I characteristic. Once forward voltage exceeds cut-in ($V_F \approx 1.8\text{--}2.2\text{ V}$), internal dynamic resistance drops to near zero. Without a series resistor $R_S = \frac{V_{CC} - V_F}{I_F}$, current escalates uncontrolled, leading to immediate thermal destruction.
\vspace{3mm}


\section{Section D: Workshop Troubleshooting \& Fault Diagnostic Case Study}

\textbf{Troubleshooting Scenario: Unresponsive Logic Probe on Trainer Module}\\
A student assembles a TTL logic circuit on the digital trainer module. However, the logic indicator LED remains permanently dark regardless of whether an input is tied to $+5\text{V}$ or $0\text{V}$.\\\\
\textbf{Diagnostic Step 1:} Check trainer DC power rails using DMM in DC voltage mode. Confirm $+5.0\text{ V} \pm 0.25\text{ V}$ between red rail and blue ground rail.\\
\textbf{Diagnostic Step 2:} Verify common ground connection. A classic workshop error is leaving the circuit ground unbonded to the trainer module ground rail, causing floating reference voltage.\\
\textbf{Diagnostic Step 3:} Check the LED polarity. The long lead (anode) must connect towards the positive logic output, with cathode to ground through the internal series limiting resistor.
\vspace{4mm}
